% ============================================================================
\clearpage
\section*{v4.9.0-r09 Source-Only: What the Local Frame Factor Can Be}
\addcontentsline{toc}{section}{v4.9.0-r09 Source-Only: What the Local Frame Factor Can Be}
\label{sec:v490-r09-local-frame-factor}
% ============================================================================

\paragraph{Purpose.}
r08 separated a fixed structural factor \(\gL=\gBI/\gstring\) from a local factor \(\gL^{(\mathrm{local})}\) that may vary.
This layer collects what the source says the local factor depends on.  It records two contradictions found on the way,
tests three readings of the local factor against clock data, and corrects the arithmetic of prediction P5.  Files:
\texttt{llm/Erkundungsprotokoll\_Weinberg\_Kegel.md} (entries E14, E15) and the script
\texttt{llm/k49\_lokaler\_faktor.py}.  No new prediction is made.  Earlier layers are not rewritten.

\noindent\textbf{Source-only status.}
This v4.9.0-r09 layer is source-only.  It does not build a PDF, does not assign a new DOI,
and does not authorize a public release.  The v4.9.0 Version DOI
\href{https://doi.org/10.5281/zenodo.23534954}{10.5281/zenodo.23534954} and the Concept DOI
\href{https://doi.org/10.5281/zenodo.19210728}{10.5281/zenodo.19210728} are unchanged by it.

\subsection*{v4.9.0-r09 what the source says}
\label{subsec:v490-r09-source}

The source names four influences on the local factor.  For none of them does it give a derived relation, and it gives
no field law for \(\gL^{(\mathrm{local})}(x,t)\).
\begin{center}
\small
\begin{tabular}{@{}>{\raggedright\arraybackslash}p{0.17\textwidth}>{\raggedright\arraybackslash}p{0.33\textwidth}>{\raggedright\arraybackslash}p{0.42\textwidth}@{}}
\toprule
influence & place & content\\
\midrule
observer & Theorem~\ref{thm:frame-offset} & every matter observer carries \(\zminus\neq0\); measured constants are frame-invariant values times a Lorentz factor of the observer\\
cosmic time & Proposition~\ref{prop:H0-tension} & \(\bsym\) grows ``as matter structures in the universe grew''\\
Earth's orbit & corollary ``\(\alpha\)-drift as a frame effect'' & the factor varies ``as Earth moves through the Sphaera-Cascade geometry over its orbit'', with a period of one year, absent for a clock that does not co-orbit\\
potential & prediction P5 near Corollary~\ref{cor:falsifiable} & amplitude of the \(\alpha\) modulation \(\sim10^{-10}\,\delta\Phi/c^2\)\\
\bottomrule
\end{tabular}
\end{center}

\paragraph{Two contradictions.}
\begin{enumerate}[leftmargin=2em]
\item Theorem~\ref{thm:frame-offset} calls \(\gBI\) ``the Lorentz-transformed \(\gstring\), observed from a matter frame''.
In that theorem \(\gBI\) depends on the observer.  Theorem~\ref{thm:k-values} sets \(k(\gBI)=0\), ``frame-invariant by
construction''.
\item Theorem~\ref{thm:frame-offset} projects all constants, including \(c\), with one Lorentz factor, which means
\(k=1\) for all.  Theorem~\ref{thm:k-values} gives \(k(c)=0\) and \(k(G_N)=2\).
\end{enumerate}
Both are resolved if Theorem~\ref{thm:frame-offset} is read as a statement about the structural factor alone.  That is the
separation of r18 and r08.  What the local factor consists of then remains open.

\subsection*{v4.9.0-r09 three readings tested}
\label{subsec:v490-r09-readings}

With the rule \(\delta\ln X=k(X)\,\delta\ln\gL^{(\mathrm{local})}\), three readings were tested.

\paragraph{Theoremlet v4.9.0-r09.1.}
\begin{enumerate}[leftmargin=2em]
\item[(V)] \emph{Velocity against a preferred frame.}  If the observer's rapidity adds to \(\eta\), Earth's orbital motion
projected on the dipole of the microwave background gives \(\delta\eta\approx9.8\times10^{-5}\).  That yields
\(\delta\ln\gL=2\tanh2\eta\,\delta\eta\approx1.6\times10^{-4}\), and with \(k(\alpha)=1\) an annual \(\alpha\) modulation of the
same size.  The clock fit quoted in r03 (F2) allows about \(2\times10^{-17}\).  The reading is \textbf{excluded}, by a factor
of about \(10^{13}\).  The static part, \(370\) km/s against the background, would make \(\alpha\) depend on direction at the
level of \(2\times10^{-3}\).
\item[(P)] \emph{Gravitational potential.}  \(\delta\ln\gL^{(\mathrm{local})}=\kappa\,\delta\Phi/c^2\).  The source's own
coefficient \(\kappa\sim10^{-10}\) and the eccentricity of Earth's orbit (\(\delta\Phi/c^2\approx1.6\times10^{-10}\)) give an
annual \(\delta\alpha/\alpha\approx1.6\times10^{-20}\), inside the bounds.  The reading \textbf{survives}.  It cannot drive the
cosmological drift of Proposition~\ref{prop:H0-tension}: potentials of large-scale structure, \(\sim10^{-5}\), give
\(\delta\ln\gL\sim10^{-15}\), against about \(5\times10^{-2}\) needed.
\item[(D)] \emph{Matter density.}  The source gives no functional form.  Density differs between laboratory, space and
astrophysical environments by many orders of magnitude, so any appreciable coupling would have to be screened.  The
reading is \textbf{not testable} as stated.
\end{enumerate}

\paragraph{Proof sketch.}
(V) \({\rm d}\ln\cosh2\eta/{\rm d}\eta=2\tanh2\eta\), with \(29.8\cos11^\circ\) km/s against \(c\).  (P) \(GM_\odot/(c^2a)=9.87\times
10^{-9}\) times the eccentricity \(0.0167\).  Checked numerically.  \(\square\)

\paragraph{Consequence.}
Of the three readings only the potential reading is compatible with the clocks.  It is also the one the source itself
states in P5.  It describes a small violation of local position invariance and cannot explain the Hubble tension.
Within it, every constant couples to the potential with \(\beta_X=k(X)\,\kappa\):
\begin{itemize}[leftmargin=2em]
\item the Compton clock of r08.2, with \(k=0\), does not couple;
\item \(\alpha\) couples with \(\kappa\);
\item an atomic clock and \(G_N\), with \(k=2\), couple with \(2\kappa\).
\end{itemize}
The ratio of the couplings, not their size, is the content of the source.  The size \(\kappa\) is not derived.

\subsection*{v4.9.0-r09 the arithmetic of P5}
\label{subsec:v490-r09-p5}

P5 states
\(\frac{\delta\alpha/\alpha}{\delta G_N/G_N}\big|_{\rm annual}=\gL^{k(\alpha)-k(G_N)}\cdot\frac{k(\alpha)}{k(G_N)}=\gL^{-1}\approx
0.537\).

\paragraph{Theoremlet v4.9.0-r09.2.}
From \(X^{(\mathrm{matter})}=\gL^{k}X^{(\mathrm{photon})}\) with fixed photon-frame values follows
\(\delta X/X=k(X)\,\delta\gL/\gL\).  Hence
\[
\frac{\delta\alpha/\alpha}{\delta G_N/G_N}=\frac{k(\alpha)}{k(G_N)}=\frac12 .
\]
The static factor \(\gL^{k(\alpha)-k(G_N)}\) does not belong in a ratio of relative changes.  The middle expression of P5
equals \(\gL^{-1}/2=0.268\), not \(\gL^{-1}\).  The corrected value \(0.5\) is the same in every reading of the local factor.

\paragraph{Proof sketch.}
Logarithmic differentiation.  \(\square\)

\subsection*{v4.9.0-r09 register entries}
\label{subsec:v490-r09-register}

\begin{center}
\small
\begin{tabular}{@{}>{\raggedright\arraybackslash}p{0.34\textwidth}>{\raggedright\arraybackslash}p{0.58\textwidth}@{}}
\toprule
statement & status after v4.9.0-r09\\
\midrule
local factor \(\gL^{(\mathrm{local})}\) & four influences named (observer, time, orbit, potential), no field law; only the potential reading survives the clocks (r09.1)\\
Thm.~\ref{thm:frame-offset} vs Thm.~\ref{thm:k-values} & two contradictions (\(\gBI\) observer-dependent vs \(k=0\); one factor for all vs \(k(c)=0\), \(k(G_N)=2\)); resolved by reading Thm.~\ref{thm:frame-offset} as structural\\
P5, annual ratio & corrected to \(k(\alpha)/k(G_N)=1/2\); the stated \(\gL^{-1}=0.537\) is wrong twice (r09.2)\\
potential coupling & \(\beta_X=k(X)\kappa\); Compton clock uncoupled; \(\kappa\sim10^{-10}\) not derived, below present bounds\\
Prop.~\ref{prop:H0-tension} & the surviving reading gives \(10^{-15}\) instead of \(5\times10^{-2}\); unchanged: hypothesis without working mechanism\\
\bottomrule
\end{tabular}
\end{center}

\subsection*{v4.9.0-r09 ledger and guard update}
\label{subsec:v490-r09-ledger-guard}

\[
R_{\rm loc}^{v490\text{-}r09}=(0,0,1)
\]
for (influences and contradictions recorded, readings of the local factor decided against clocks, size of \(\kappa\) open).

The layer spends no guard:
\[
R_{\rm glob}^{(v490\text{-}r09,+)}=(0,1,1,1,1),
\qquad
\Gamma_{\rm guard}^{v490\text{-}r09}=(1,1,1,1,1),
\qquad
\beta_4=1/256\ \text{unspent}.
\]
It does not imply exact finite \(C5\) zero closure, exact all-readout terminality,
\(C6/T2\), a global photon-clock/RH theorem, or global all-readout closure.

